\documentclass[a4paper,12pt]{article}
\usepackage[utf8]{inputenc}
\usepackage[T1]{fontenc}
\usepackage[french]{babel}
\usepackage{geometry}
\usepackage{xcolor}
\usepackage[explicit]{titlesec}
\usepackage{tcolorbox}
\usepackage{lmodern}
\usepackage{enumitem}
\usepackage{booktabs}
\usepackage{tabularx}
\usepackage{array}
\usepackage{fancyhdr}
\usepackage{hyperref}
\usepackage{graphicx}
\usepackage{float}

\geometry{margin=2.5cm, top=3cm, bottom=3cm}

% ── Couleurs ──────────────────────────────────────────────────────────────────
\definecolor{primary}{RGB}{0, 70, 140}
\definecolor{secondary}{RGB}{230, 240, 255}
\definecolor{lightgray}{RGB}{245, 245, 245}
\definecolor{darkgray}{RGB}{80, 80, 80}

% ── En-tête / Pied de page ────────────────────────────────────────────────────
\pagestyle{fancy}
\fancyhf{}
\lhead{\textcolor{darkgray}{\small BTS SIO — SISR | AP-2}}
\rhead{\textcolor{darkgray}{\small Projet DYLAN — M2L}}
\cfoot{\textcolor{darkgray}{\thepage}}
\renewcommand{\headrulewidth}{0.4pt}

% ── Style des sections ────────────────────────────────────────────────────────
\titleformat{\section}
  {\normalfont\Large\bfseries\sffamily\color{primary}}
  {}
  {0pt}
  {\colorbox{secondary}{\parbox{\dimexpr\textwidth-2\fboxsep}{\thesection\quad #1}}}

\titleformat{\subsection}
  {\normalfont\normalsize\bfseries\sffamily\color{primary}}
  {\thesubsection}
  {0.5em}
  {#1}
  [\vspace{-0.3em}\textcolor{primary!40}{\hrule height 0.5pt}\vspace{0.3em}]

% ── Style des listes ──────────────────────────────────────────────────────────
\setlist[itemize,1]{label=\textcolor{primary}{$\bullet$}, leftmargin=2em, labelsep=1em}
\setlist[itemize,2]{label=\textcolor{primary!60}{--}, leftmargin=2em, labelsep=0.8em}
\setlist[enumerate,1]{label=\textcolor{primary}{\arabic*.}, leftmargin=2em, labelsep=1em}
\setlist{itemsep=2pt, parsep=2pt}

% ── Boîte info ────────────────────────────────────────────────────────────────
\tcbuselibrary{skins}
\newtcolorbox{infobox}{
  colback=secondary, colframe=primary,
  boxrule=0.5pt, leftrule=4pt,
  sharp corners, left=8pt, right=8pt, top=6pt, bottom=6pt
}

% ── Tableau stylé ─────────────────────────────────────────────────────────────
\newcolumntype{L}[1]{>{\raggedright\arraybackslash}p{#1}}
\newcolumntype{C}[1]{>{\centering\arraybackslash}p{#1}}

\hypersetup{colorlinks=true, linkcolor=primary, urlcolor=primary, citecolor=primary}

% =============================================================================
\begin{document}
\thispagestyle{empty}

% ── Page de garde ─────────────────────────────────────────────────────────────
\begin{center}
  \vspace*{2cm}
  {\color{primary}\rule{\textwidth}{2pt}}
  \vspace{0.8em}

  {\Huge\bfseries\sffamily\color{primary} Projet DYLAN}\\[0.4em]
  {\Large\sffamily\color{darkgray} Insertion dynamique dans le VLAN par adresse MAC}\\[0.2em]
  {\normalsize\color{darkgray} Maison des Ligues de Lorraine — AP-2 SISR}

  \vspace{0.8em}
  {\color{primary}\rule{\textwidth}{2pt}}

  \vspace{2cm}

  \begin{tcolorbox}[colback=lightgray, colframe=primary!30, sharp corners, width=0.7\textwidth]
    \centering
    \textbf{Groupe :} Maël — Joan — Camille — Anna\\[0.3em]
    \textbf{Formation :} BTS SIO option SISR\\[0.3em]
    \textbf{Année :} 2025–2026\\[0.3em]
    \textbf{Date :} \today
  \end{tcolorbox}
\end{center}

\vfill
{\small\color{darkgray}\centering Réseau CERTA — Contexte M2L\par}

\newpage
\tableofcontents
\newpage

% =============================================================================
\section{Contexte et problématique}

\subsection{Présentation du contexte}

% Depuis la mise en place des VLANs (projet VELANNE), les personnels des ligues
% ne peuvent plus accéder à leurs données depuis la salle de réunion commune.
% Les postes fixes sont dans leur VLAN de ligue, mais la prise de la salle de
% réunion est partagée entre plusieurs ligues.

\subsection{Problématique identifiée}

% Comment permettre à un portable de se placer automatiquement
% dans le bon VLAN selon la ligue à laquelle il appartient,
% sans intervention manuelle sur le switch ?

\subsection{Objectifs du projet}

\begin{itemize}
  \item
  \item
  \item
\end{itemize}

% =============================================================================
\section{Infrastructure réseau}

\subsection{Matériel utilisé}

\begin{itemize}
  \item \textbf{Switch Cisco SF-300} — Switch administrable supportant la gestion de VLANs par adresse MAC (bâtiments A et B)
  \item \textbf{PC fixes des ligues} — Connectés sur les ports ACCESS de leur VLAN respectif
  \item \textbf{Portables des ligues} — Se connectent sur la prise de la salle de réunion
\end{itemize}

\subsection{Schéma réseau}

% Insérer ici le schéma de topologie : switch SF-300, ports bureaux (ACCESS),
% port salle de réunion (GÉNÉRAL), postes fixes et portables

\begin{infobox}
  \textit{Schéma à insérer ici — topologie switch Cisco SF-300 avec les VLANs}
\end{infobox}

\subsection{Plan de VLANs}

\begin{infobox}
  \textbf{Convention de numérotation M2L :} \texttt{[Bâtiment][Étage][Numéro ligue]}\\
  Exemple : VLAN \textbf{121} = Bâtiment A, Étage 2, Ligue 1 (Volley)
\end{infobox}

\vspace{0.5em}
\begin{tabularx}{\textwidth}{C{2cm} C{3cm} C{3cm} X}
  \toprule
  \textbf{VLAN ID} & \textbf{Ligue} & \textbf{Groupe MAC} & \textbf{Équipements concernés} \\
  \midrule
  121 & Volley & Groupe-121 & Postes fixes volley + portables volley \\
  122 & Escrime & Groupe-122 & Postes fixes escrime + portables escrime \\
  123 & Badminton & Groupe-123 & Postes fixes badminton + portables badminton \\
  \bottomrule
\end{tabularx}

\subsection{Adresses MAC des équipements}

\begin{tabularx}{\textwidth}{L{4cm} C{3cm} C{3cm} C{3cm}}
  \toprule
  \textbf{Équipement} & \textbf{Adresse MAC} & \textbf{Groupe} & \textbf{VLAN} \\
  \midrule
  Portable Volley & & Groupe-121 & 121 \\
  Portable Escrime & & Groupe-122 & 122 \\
  Portable Badminton & & Groupe-123 & 123 \\
  \bottomrule
\end{tabularx}

% =============================================================================
\section{Configuration du switch Cisco SF-300}

\subsection{Création des VLANs}

% Aller dans VLAN Management > VLAN Settings > Add
% Créer VLAN 121, 122, 123

\begin{tcolorbox}[colback=black!90, coltext=green!70!white, fontupper=\ttfamily\small,
  title={\color{white}\small Création des VLANs (CLI)}, colbacktitle=primary]
% vlan database
% vlan 121 name Volley
% vlan 122 name Escrime
% vlan 123 name Badminton
% exit
\end{tcolorbox}

\subsection{Configuration des ports en mode ACCESS}

% Chaque port de bureau est assigné au VLAN de la ligue correspondante.
% Un port en mode ACCESS ne transmet que les trames d'un seul VLAN.

\begin{tcolorbox}[colback=black!90, coltext=green!70!white, fontupper=\ttfamily\small,
  title={\color{white}\small Ports ACCESS (CLI)}, colbacktitle=primary]
% interface GigabitEthernet X
%  switchport mode access
%  switchport access vlan 121   (adapter selon la ligue)
% exit
\end{tcolorbox}

\begin{tabularx}{\textwidth}{C{3cm} C{3cm} C{3cm} X}
  \toprule
  \textbf{Port} & \textbf{Mode} & \textbf{VLAN} & \textbf{Équipement connecté} \\
  \midrule
   & ACCESS & 121 & Poste fixe Volley \\
   & ACCESS & 122 & Poste fixe Escrime \\
   & ACCESS & 123 & Poste fixe Badminton \\
  \bottomrule
\end{tabularx}

\subsection{Configuration du port de la salle de réunion (mode GÉNÉRAL)}

% Le port de la prise de la salle de réunion est mis en mode GÉNÉRAL :
% il est membre des VLANs 121, 122 et 123 simultanément.
% C'est sur ce port que le mécanisme MAC-based VLAN s'applique.

\begin{tcolorbox}[colback=black!90, coltext=green!70!white, fontupper=\ttfamily\small,
  title={\color{white}\small Port salle de réunion — mode GÉNÉRAL (CLI)}, colbacktitle=primary]
% interface GigabitEthernet X
%  switchport mode general
%  switchport general allowed vlan add 121,122,123 untagged
% exit
\end{tcolorbox}

\subsection{Création des groupes d'adresses MAC}

% Via l'interface web : VLAN Management > MAC-Based Groups > Add
% Créer un groupe par ligue et y ajouter les adresses MAC des portables

\begin{infobox}
  \textbf{Chemin interface web :} \texttt{VLAN Management > MAC-Based Groups > Add}\\
  Créer un groupe par ligue (Groupe-121, Groupe-122, Groupe-123) et y ajouter
  les adresses MAC des portables correspondants.
\end{infobox}

\begin{tcolorbox}[colback=black!90, coltext=green!70!white, fontupper=\ttfamily\small,
  title={\color{white}\small Groupes MAC (CLI)}, colbacktitle=primary]
% mac address-table static XX:XX:XX:XX:XX:XX vlan 121 interface GigabitEthernet X
% (répéter pour chaque adresse MAC)
\end{tcolorbox}

\subsection{Mapping Groupe MAC → VLAN}

% Via l'interface web : VLAN Management > MAC-Based Groups to VLAN > Add
% Associer chaque groupe à son VLAN cible

\begin{infobox}
  \textbf{Chemin interface web :} \texttt{VLAN Management > MAC-Based Groups to VLAN > Add}\\
  Associer : Groupe-121 → VLAN 121, Groupe-122 → VLAN 122, Groupe-123 → VLAN 123.
\end{infobox}

\begin{tabularx}{\textwidth}{C{3.5cm} C{3.5cm} X}
  \toprule
  \textbf{Groupe MAC} & \textbf{VLAN cible} & \textbf{Description} \\
  \midrule
  Groupe-121 & 121 & Portables de la ligue Volley \\
  Groupe-122 & 122 & Portables de la ligue Escrime \\
  Groupe-123 & 123 & Portables de la ligue Badminton \\
  \bottomrule
\end{tabularx}

% =============================================================================
\section{Tests et validation}

\subsection{Plan de tests}

\begin{tabularx}{\textwidth}{C{0.5cm} L{4.5cm} L{4cm} C{2cm} C{2cm}}
  \toprule
  \textbf{\#} & \textbf{Test} & \textbf{Action} & \textbf{Résultat attendu} & \textbf{Résultat obtenu} \\
  \midrule
  1 & Affectation automatique VLAN 121 & Brancher le portable Volley sur la prise salle de réunion & VLAN 121 assigné & \\
  2 & Affectation automatique VLAN 122 & Brancher le portable Escrime & VLAN 122 assigné & \\
  3 & Affectation automatique VLAN 123 & Brancher le portable Badminton & VLAN 123 assigné & \\
  4 & Communication dans le même VLAN & Ping poste fixe Volley depuis portable Volley & Succès & \\
  5 & Isolation entre VLANs & Ping poste fixe Escrime depuis portable Volley & Échec (isolé) & \\
  6 & Partage de dossiers Windows & Accéder au dossier partagé depuis le portable dans le même VLAN & Succès & \\
  \bottomrule
\end{tabularx}

\subsection{Captures d'écran}

\begin{infobox}
  \textit{Captures d'écran à insérer ici : interface SF-300, assignation VLAN, résultats de ping, partage réseau Windows}
\end{infobox}

% =============================================================================
\section{Avantages et inconvénients de la solution}

\begin{tabularx}{\textwidth}{X X}
  \toprule
  \textbf{Avantages} & \textbf{Inconvénients} \\
  \midrule
  & \\
  & \\
  & \\
  & \\
  \bottomrule
\end{tabularx}

% =============================================================================
\section{Support de formation — Partage de dossiers Windows}

\subsection{Objectif}

% Expliquer comment deux postes dans le même VLAN peuvent partager des dossiers
% sous Windows (partage réseau, droits d'accès, accès par \\NomMachine\Dossier)

\subsection{Prérequis}

\begin{itemize}
  \item Les deux machines doivent être dans le même VLAN (et donc sur le même sous-réseau)
  \item La découverte réseau et le partage de fichiers doivent être activés dans Windows
  \item
\end{itemize}

\subsection{Procédure}

\begin{enumerate}
  \item Activer le partage réseau sur Windows :
  \begin{itemize}
    \item Panneau de configuration → Centre réseau et partage → Modifier les paramètres de partage avancés
    \item Activer la découverte réseau et le partage de fichiers
  \end{itemize}
  \item Partager un dossier :
  \begin{itemize}
    \item Clic droit sur le dossier → Propriétés → Partage → Partager
    \item Définir les droits d'accès (lecture / lecture-écriture)
  \end{itemize}
  \item Accéder au dossier partagé depuis l'autre machine :
  \begin{itemize}
    \item Explorateur de fichiers → \texttt{\textbackslash\textbackslash NomMachine\textbackslash NomDossier}
    \item Ou via l'adresse IP : \texttt{\textbackslash\textbackslash 10.x.x.x\textbackslash NomDossier}
  \end{itemize}
\end{enumerate}

\subsection{Captures d'écran}

\begin{infobox}
  \textit{Captures d'écran à insérer ici : configuration partage, accès depuis l'autre poste}
\end{infobox}

% =============================================================================
\section*{Conclusion}
\addcontentsline{toc}{section}{Conclusion}

% Résumer ce qui a été mis en place, les compétences mobilisées (SISR2),
% et l'intérêt de la solution MAC-based VLAN dans un contexte professionnel réel.

\end{document}
